% PACCHETTI FONDAMENTLAI
\usepackage[T1]{fontenc}
\usepackage[utf8]{inputenc}
\usepackage[italian]{babel}
\usepackage[babel]{csquotes}
\usepackage[style=numeric, hyperref]{biblatex}
\usepackage{microtype}
\usepackage{lmodern}
\usepackage{graphicx} % inserire immagini
\usepackage{subfig} % sottofigure
\usepackage{multicol} % due colonne
% \usepackage{ulem} % sottolineare
\usepackage{lipsum} % lorem ipsum
\usepackage{xcolor} % colori in latex
\usepackage{parskip} % rimuove l'indentazione dei nuovi paragrafi %% Add parbox=false to all new tcolorbox
\usepackage{centernot}
\usepackage[outline]{contour}\contourlength{3pt}
\usepackage{fancyhdr}
\usepackage{layout}
\usepackage[intlimits]{empheq} % Riquadri colorati attorno alle equazioni
\usepackage[most]{tcolorbox} % Riquadri colorati
\usepackage{ifthen} % IFTHEN
\usepackage{rotating} %figure ruotate 

% pacchetti matematica
\usepackage[intlimits]{amsmath} 
\usepackage{amssymb}
\usepackage{amsthm}
\usepackage{yhmath}
\usepackage{dsfont}
\usepackage{mathrsfs}
\usepackage{cancel} % semplificare
\usepackage{polynom} %divisione tra polinomi
\usepackage{forest} % grafi ad albero
\usepackage{booktabs} % tabelle
\usepackage{commath} %simboli e differenziali
\usepackage{bm} %bold
\usepackage[fulladjust]{marginnote} %to use marginnote for date notes
\usepackage{arrayjobx}%array
\usepackage{mathtools}
%%%%%%%%%%%%%


%%%% CONTATORE DIMOSTRAZIONI
\newcounter{numerodidimostrazioni}
\newcommand{\mostranumerodimostrazioni}{\cleardoublepage\thispagestyle{empty}\begin{center}\null\vfill \textit{Questo corso comprende {\thenumerodidimostrazioni} dimostrazioni}\vfill\null\end{center}}
%%%%

%%%% QUIVER
\newcommand{\duepunti}{\,\mathchar\numexpr"6000+`:\relax\,}
% A TikZ style for curved arrows of a fixed height, due to AndréC.
\tikzset{curve/.style={settings={#1},to path={(\tikztostart)
    .. controls ($(\tikztostart)!\pv{pos}!(\tikztotarget)!\pv{height}!270:(\tikztotarget)$)
    and ($(\tikztostart)!1-\pv{pos}!(\tikztotarget)!\pv{height}!270:(\tikztotarget)$)
    .. (\tikztotarget)\tikztonodes}},
    settings/.code={\tikzset{quiver/.cd,#1}
        \def\pv##1{\pgfkeysvalueof{/tikz/quiver/##1}}},
    quiver/.cd,pos/.initial=0.35,height/.initial=0}

% TikZ arrowhead/tail styles.
\tikzset{tail reversed/.code={\pgfsetarrowsstart{tikzcd to}}}
\tikzset{2tail/.code={\pgfsetarrowsstart{Implies[reversed]}}}
\tikzset{2tail reversed/.code={\pgfsetarrowsstart{Implies}}}
% TikZ arrow styles.
\tikzset{no body/.style={/tikz/dash pattern=on 0 off 1mm}}
%%%%%%%%%%

%%% HEADER e FOOTER
%% HEAD and FOOT
\newcommand{\datainizion}{testo}
\newcommand{\hdrnew}[1]{
	\fancyhead{} % cancella tutti i campi
	\fancyfoot{}
	\fancyhead[LE]{#1}
	\fancyhead[RO]{\nouppercase{\leftmark}}
	% \fancyfoot[LE,RO]{\thepage}
	\fancyfoot[C]{\thepage}
	% \fancyfoot[CO,RE]{Per: Dean A. Smith}
	\renewcommand{\headrulewidth}{0.4pt}
	\renewcommand{\footrulewidth}{0pt}
	\pagestyle{fancy}
}
\newcommand{\days}[1]{\renewcommand{\datainizion}{#1}}
%%%%%%%%

%% DEFINIZIONI COMANDI MATEMATICI
\let\sin\relax %TOGLIE LA DEFINIZIONE SU "\sin"

% cambia la definizione di empty set
% ---
\let\oldemptyset\emptyset
% ---
\let\emptyset\varnothing 
% ---
% \let\emptyset\relax
% \newcommand{\emptyset}{\text{\textnormal{\O}}}
% ---

\DeclareMathOperator{\bounded}{bd}
\DeclareMathOperator{\sin}{sen}
\DeclareMathOperator{\epi}{Epi}
\DeclareMathOperator{\cl}{cl}
\DeclareMathOperator{\graph}{graph}
\DeclareMathOperator{\arcsec}{arcsec}
\DeclareMathOperator{\arccot}{arccot}
\DeclareMathOperator{\arccsc}{arccsc}
\DeclareMathOperator{\spettro}{Spettro}
\DeclareMathOperator{\nulls}{nullspace}
\DeclareMathOperator{\dom}{dom}
\DeclareMathOperator{\ar}{ar}
\DeclareMathOperator{\const}{Const}
\DeclareMathOperator{\fun}{Fun}
\DeclareMathOperator{\rel}{Rel}
\DeclareMathOperator{\altezza}{ht}
\let\det\relax %TOGLIE LA DEFINIZIONE SU "\det"
\DeclareMathOperator{\det}{det}
\DeclareMathOperator{\End}{End}
\DeclareMathOperator{\gl}{GL}
\DeclareMathOperator{\Id}{Id}
\DeclareMathOperator{\id}{Id}
\DeclareMathOperator{\I}{\mathds{1}}
\DeclareMathOperator{\II}{II}
\DeclareMathOperator{\rank}{rank}
\DeclareMathOperator{\tr}{tr}
\DeclareMathOperator{\tc}{t.c.}
\DeclareMathOperator{\T}{T}
\DeclareMathOperator{\var}{Var}
\DeclareMathOperator{\cov}{Cov}
\DeclareMathOperator{\st}{st}
\DeclareMathOperator{\mon}{Mon}
\newcommand{\trasposta}[1]{\prescript{\text{T}}{}{#1}}
\newcommand{\nonstandard}[1]{\prescript{*}{}{#1}}
\newcommand{\halo}{\mathfrak{h}}
\newcommand{\hhalo}[1]{\prescript{\halo}{}{#1}}
%%% Modo 1

% \newcommand{\sequence}[2][n \in \N]{{\langle #2\rangle}^{#1}
% }
% \newcommand{\filtersequence}[2][n \in \N]{\sequence[#1]{#2}_{\mathcal{U}}}

%%%
%%% Modo 2

\newcommand{\sequence}[2][n \in \N]{{\langle #2\rangle}_{\!#1}
}
\newcommand{\filtersequence}[2][n \in \N]{%
	{\overline{\langle #2\rangle}}_{\!#1}%
}

%%%

%%% Modo 3

% \newcommand{\sequence}[2][n \in \N]{{\langle #2\rangle}_{#1}
% }
% \newcommand{\filtersequence}[2][n \in \N]{\sequence[#1]{#2}^{\mathcal{U}}}

%%%
\newcommand{\R}{\mathds{R}}
\newcommand{\starR}{\nonstandard{\R}}
\newcommand{\borel}{\mathscr{B}}
\newcommand{\lebesgue}[1]{\mathscr{L}\left(#1\right)}
\newcommand{\media}{\mathds{E}}
\newcommand{\K}{\mathds{K}}
\newcommand{\A}{\mathds{A}}
\newcommand{\Q}{\mathds{Q}}
\newcommand{\N}{\mathds{N}}
\newcommand{\C}{\mathds{C}}
\newcommand{\Z}{\mathds{Z}}
\newcommand{\rmn}{\R^{m,n}}
\newcommand{\qo}{\hspace{1em}\text{q.o.}\,}
\newcommand{\pprec}{\mathrel{\prec\! \prec}}
\newcommand{\npprec}{\mathrel{\prec\! \nprec}}
\renewcommand{\tilde}[1]{\widetilde{#1}}
\renewcommand{\parallel}{\mathrel{/\mkern-5mu/}}
\newcommand{\cardin}[1]{\left\vert #1\right\vert}
\newcommand{\parti}[1]{\wp (#1)}
\newcommand{\diff}[1]{\operatorname{d}_{#1}}
\renewcommand{\vec}[1]{\overrightarrow{\vphantom{i}#1}}
\newcommand{\floor}[1]{\left\lfloor #1 \right\rfloor}

%tramite i prossimi due comandi posso decidere come scrivere i logaritmi naturali in tutti i documenti: ho infatti eliminato qualsiasi differenza tra "ln" e "log": se si vuole qualcosa di diverso bisogna inserire manualmente il tutto
\let\ln\relax
\DeclareMathOperator{\ln}{ln}
\let\log\relax
\DeclareMathOperator{\log}{log}
%%%%%%

%% NUOVI COMANDI
\newcommand{\straniero}[1]{\textit{#1}} %parole straniere
\newcommand{\titolo}[1]{\textsc{#1}} %titoli
\newcommand{\qedd}{\tag*{$\blacksquare$}} %qed per ambienti matemastici
\renewcommand{\qedsymbol}{$\blacksquare$} %modifica colore qed
\newcommand{\ooverline}[1]{\overline{\overline{#1}}}
\newcommand{\circoletto}[1]{\left(#1\right)^\text{o}}
%
\newcommand{\qmatrice}[1]{\begin{pmatrix}
#1_{11} & \cdots & #1_{1n}\\
\vdots & \ddots & \vdots \\
#1_{m1} & \cdots & #1_{mn}
\end{pmatrix}}
%
\newcommand{\parentesi}[2]{%
\underset{#1}{\underbrace{#2}}%
}
%
\newcommand{\norma}[1]{% Norma 
\left\lVert#1\right\rVert%
}
\newcommand{\scalare}[2]{% Scalare
\left\langle #1, #2\right\rangle
}
%%%%%

%% RESTRIZIONI
\newcommand{\referenze}[2]{
	\phantomsection{}#2\textsuperscript{\textcolor{blue}{\textbf{#1}}}
}
\def\restriction#1#2{\mathchoice
              {\setbox1\hbox{${\displaystyle #1}_{\scriptstyle #2}$}
              \restrictionaux{#1}{#2}}
              {\setbox1\hbox{${\textstyle #1}_{\scriptstyle #2}$}
              \restrictionaux{#1}{#2}}
              {\setbox1\hbox{${\scriptstyle #1}_{\scriptscriptstyle #2}$}
              \restrictionaux{#1}{#2}}
              {\setbox1\hbox{${\scriptscriptstyle #1}_{\scriptscriptstyle #2}$}
              \restrictionaux{#1}{#2}}}
\def\restrictionaux#1#2{{#1\,\smash{\vrule height .8\ht1 depth .85\dp1}}_{\,#2}} 
%%%%%%%%%%%

%% SEZIONE GRAFICA
\usepackage{tikz}
\usetikzlibrary{matrix, patterns, calc, decorations.pathreplacing, hobby, decorations.markings, decorations.pathmorphing, babel}
\usepackage{tikz-3dplot}
\usepackage{mathrsfs} %per geogebra
\usepackage{tikz-cd}
\tikzset
{
  %surface/.style={fill=black!10, shading=ball,fill opacity=0.4},
  plane/.style={black,pattern=north east lines},
  curve/.style={black,line width=0.5mm},
  dritto/.style={decoration={markings,mark=at position 0.5 with {\arrow{Stealth}}}, postaction=decorate},
  rovescio/.style={decoration={markings,mark=at position 0.5 with {\arrow{Stealth[reversed]}}}, postaction=decorate}
}
\usepackage{pgfplots} % stampare le funzioni
	\pgfplotsset{/pgf/number format/use comma,compat=1.15}
	%\pgfplotsset{compat=1.15} %per geogebra
	\usepgfplotslibrary{fillbetween, polar}
%%%%%%

%% CITAZIONI
\usepackage{lineno}

\newcommand{\citazione}[1]{%
  \begin{quotation}
  \begin{linenumbers}
  \modulolinenumbers[5]
  \begingroup
  \setlength{\parindent}{0cm}
  \noindent #1
  \endgroup
  \end{linenumbers}
  \end{quotation}\setcounter{linenumber}{1}
  }
%%%%%%


\usepackage{hyperref}
\hypersetup{%
	pdfauthor={Davide Peccioli},
	pdfsubject={Appunti UniTO},
	allcolors=black,
	citecolor=black,
	colorlinks=true, 
	bookmarksopen=true}
